\chapter*{Preface}
\addcontentsline{toc}{chapter}{Preface}

This is a successor volume to \textit{Proof and Countermodel: An
Introduction to Logic with Lean}. It assumes that book's method —
every claim proved where it holds, paired with an explicit
countermodel where a plausible-looking converse or strengthening
fails — but not its content. A reader who has completed the
introductory sequence needs no further preparation; a reader who
has not will find the method explained here only briefly, since it
was that earlier book's whole subject.

Where \textit{Proof and Countermodel} stayed classical throughout
and stopped at elementary set theory, this book goes three places
that book deliberately left aside: logics in which excluded middle
or explosion fail, modal structures in which truth depends on which
state one is reasoning from, and — the point the whole volume is
organized around — the idea that admissibility, repair, and
continuation are themselves best understood as a logic, with their
own consequence relation, rather than as metaphors borrowed from
elsewhere. A closing chapter sketches metatheory: soundness for the
fragments developed, decidability where it holds, and a small
demonstration of syntax represented and inspected from inside a
formal system.

The title names the book's actual concern more directly than
\textit{Proof and Countermodel} needed to: what follows from what,
once the space of admissible inferences is itself allowed to be
governed by constraints that classical logic does not impose.
